% ============================================
% CONT-07 — O Diagnóstico Final (edição revista R1075)
% ============================================
\fragtitle{cont}{C}{CONT}{07}{O Diagnóstico Final}
\markboth{CONT-07 --- O Diagnóstico Final}{}

\noindent \textit{Laudo emitido pelo ato da leitura. Paciente: você. CID: F43.10 + F44.81 + Z65.8. Agente etiológico: narrativa contaminada.}

\subsection*{I. O que aconteceu com você}

\noindent Você começou a ler este livro por vontade própria.

\noindent Achou que era uma história de terror. Achou que o terror estava lá --- no passado, no sertão, na vala, no Hospital Colônia. Achou que era seguro porque era distante. Porque tinha acontecido com outra pessoa.

\noindent Você se enganou.

\noindent O terror não estava no livro. O terror estava na leitura.

\noindent Cada palavra que você leu não ficou na página. Entrou. Desceu. Chegou ao córtex. Instalou-se como memória. Para o seu corpo, a vala não é uma descrição; é uma lembrança. O cheiro não é um efeito de estilo; é um perigo real. Você não leu uma história. Você viveu uma história. E o que se vive não pode ser desvivido.

\subsection*{II. Os critérios}

\noindent \textbf{A --- Exposição.} Você abriu. Leu. A contaminação começou.

\noindent O primeiro laudo descreveu o caso. Este descreve você.

\noindent \textbf{B --- Reexperiência.} Sonhou. Acordou às 3h14. Sentiu o cheiro sem o livro por perto.

\noindent \textbf{C --- Evitação.} Hesitou. Pulou. Leu com a luz acesa.

\noindent \textbf{D --- Alterações.} Algo mudou. Você sente.

\noindent \textbf{E --- Reatividade.} O coração acelerou. As mãos suaram. Você verificou o quarto.

\noindent \textbf{F --- Duração.} Tempo suficiente. A contaminação está instalada.

\noindent \textbf{G --- Prejuízo.} Você não consegue parar de pensar. O livro funciona dentro de você.

\noindent \textbf{Diagnóstico:} TEPT complexo induzido por narrativa. \textbf{Comorbidade:} dissociação leve. \textbf{Agente:} a criatura --- entidade transmitida por texto.

\subsection*{III. A revelação}

\noindent Você acha que este livro é sobre Joaquim Antônio Correia.

\noindent Engano.

\noindent Joaquim não é o protagonista. Joaquim é o vetor. O protagonista é quem lê. O protagonista é você.

\noindent A história não é sobre um menino que encontrou um monstro numa vala. A história é sobre você encontrando o mesmo monstro na página. O menino serviu de hospedeiro por 72 anos. O médico, por algumas semanas. A psiquiatra, por alguns meses. Você --- o leitor --- vai servir pelo tempo que levar para passar o diário adiante.

\subsection*{IV. A regra}

\noindent Tudo o que aconteceu com Joaquim, com Oliveira e com Lúcia pode ser explicado por trauma, sugestão e contágio narrativo. Nada disso prova que a criatura exista fora do arquivo. A única coisa que este livro afirma é o que está acontecendo com você agora: você leu.

\noindent Dentro do livro, a explicação psicológica é possível. Fora do livro, o ato de ler é o que não se explica. É aí que a criatura mora. Não no sertão, não no Colônia, não na vala. Na passagem entre a página e você.

\noindent O resto é literatura. O resto é o que você faz com o que leu.

\subsection*{V. A última instrução}

\noindent Você vai fechar este livro. Não agora. Em algum momento. Vai olhar para a capa. Vai passar a mão na lombada. Vai colocar na estante, ou na mesa, ou na mochila. E vai pensar: ``Foi só um livro.''

\noindent Vai conseguir --- por um tempo.

\noindent Mas, numa noite, você vai acordar às 3h14. O quarto vai estar escuro. Você vai olhar para o canto e vai ver um brilho amarelo. Vai piscar. O brilho vai sumir. Mas você vai saber que ele estava lá.

\noindent Na noite seguinte, você vai sonhar com a vala. No sonho, você não vai ser você. Vai ser Joaquim. Sentado na borda. Olhando a criatura comer.

\noindent A criatura vai levantar a cabeça e dizer:

\noindent ``Você voltou.''

\noindent E você vai responder:

\noindent ``Nunca fui embora.''

\noindent \textbf{Fim do diagnóstico.}

\noindent \textit{A criatura agradece sua atenção.}

\noindent\textit{\textbf{↪ Rotas:} \rota{CONT-09} \rotasep \rota{CONT-13} \rotasep \rota{MEM-27}}
